\section{先前工作与当前趋势 (Previous Work and Present Trends)}\label{sec:previouswork}

在以太坊\emph{黄皮书}首次发表以来的几年中，区块链开发领域得到了极大的发展。除了可扩展性之外，开发工作还围绕底层共识算法、智能合约语言和机器以及整体状态环境展开。虽然这些主题很有趣，但它们大多超出了本文的范围，因为它们通常不影响底层可扩展性。

\subsection{Polkadot}

为了提供服务，JAM 借用了与 Polkadot 相同的大量博弈论和密码学机制，称为 \textsc{Elves}，由 \cite{cryptoeprint:2024/961} 描述。然而，在 JAM 提供的实际服务方面存在重大差异，它提供的抽象更接近于其经济激励的验证者节点所产生的实际计算模型。

原始 Polkadot 提案（一个可扩展的异构多链）的一个主要要点是通过在多台主机上分区和分配工作负载来实现高性能。在这样做时，它明确表明可组合性将会降低。Polkadot 的组成部分——平行链——实际上在其性质上是高度隔离的。虽然存在一个消息传递系统（\textsc{xcmp}），但它是异步的、粗粒度的，并且由于依赖于一种高层的、缓慢演化的交互语言 \textsc{xcm} 而在实际中受到限制。

因此，Polkadot 在其组成链之间提供的可组合性低于以太坊类智能合约系统——后者提供单一且通用的对象环境，并允许支撑其成功的那种敏捷和创新的集成。就目前而言，Polkadot 是一个独立生态系统的集合，只有有限的协作机会，在人体工学上与桥接区块链非常相似，尽管具有截然不同的安全特征。一个名为 \textsc{spree} 的技术提案将利用 Polkadot 独特的共享安全性来改善可组合性，但区块链仍将保持隔离状态。

实施和推出一条区块链是困难的、耗时的且成本高昂的。根据其原始设计，Polkadot 将能够使用其服务的客户端限制为那些既有能力做到这一点、又能筹集足够保证金来赢得长期插槽拍卖的客户端，目前大约有 50 个插槽。虽然本身不是许可制的，但可访问性与以太坊等智能合约系统相比显著降低。

让尽可能多的创新者参与并相互互动，以及与彼此的用户群互动，似乎是 Web3 应用平台成功的重要组成部分。因此，可访问性至关重要。

\subsection{Ethereum}

以太坊协议在本文的精神前身——\emph{黄皮书}中被正式定义，由 \cite{wood2014ethereum} 完成。这在很大程度上源自 \cite{buterin2013ethereum} 的初始概念论文。自\emph{黄皮书}发表以来的十年中，\emph{事实上的}以太坊协议和公共网络实例经历了多次演化，主要围绕通过交易格式以及其脚本核心——以太坊虚拟机（\textsc{evm}）——的指令集和"预编译"（小众的、复杂的附加指令）来引入灵活性。

近一百万个加密经济参与者参与以太坊的验证。\footnote{实际问题确实限制了真正去中心化的程度。验证者软件明确提供了允许单个实例配置多个密钥集的功能，系统性地促成了远低于表面参与者数量的实际去中心化水平，无论是从单个运营商还是硬件角度来看。根据 \cite{hildobby2024eth2} 关于以太坊 2 的数据汇总，可以看到一个主要节点运营商 Lido 稳定地占据了近一百万加密经济参与者的近三分之一。}区块扩展通过随机化的领导者轮换方法完成，其中领导者的物理地址在其区块生产之前是公开的。\footnote{以太坊的开发者希望将其改为更安全的方案，但没有确定的时间表。}以太坊使用 \cite{buterin2019casper} 引入的 Casper-\textsc{ffg} 来确定最终性，凭借庞大的验证者基数，大约每 13 分钟完成一次链扩展的最终确认。

以太坊的直接计算性能与 2015 年推出时大致相似，一个显著的例外是现在额外的服务允许每个区块托管 1\textsc{mb} 的\emph{承诺数据}（所有节点在有限时间内存储它）。该数据不能被主状态转换函数直接使用，但特殊函数提供了数据（或其某个子部分）可用性的证明。根据 \cite{ethereum2024danksharding}，当前的设计方向是在未来几年中通过一种名为\emph{Dank-sharding}的协议，在验证者基数之间分配其存储责任来改进这一点。

根据 \cite{ethereum2024sigital}，以太坊的扩展策略是将这种数据可用性与\emph{roll-up}的私人市场相结合——各种设计的侧带计算设施，其中基于 \textsc{zk-snark} 的 roll-up 是首选。每个供应商的 roll-up 设计、执行和运营都有其自身的含义。

人们可能合理地假设，通过多供应商 roll-up 的多元化市场扩展方法将允许设计良好的解决方案蓬勃发展。然而，该策略面临潜在问题。 \cite{sharma2024ethereums} 关于各种 roll-up 去中心化程度的研究报告发现了广泛的中心化模式，但指出正在进行的工作试图缓解这一问题。它们最终能达到何种去中心化程度仍有待观察。

各种竞争供应商组成的大拼盘 roll-up 的主要区域之间，异构的通信特性（如数据报延迟和语义范围）、安全特性（如回退、损坏、停机和审查的成本）和经济特性（接受和处理某些传入消息或交易的成本）可能存在差异，甚至可能相当巨大。虽然整个以太坊网络最终可能提供执行侧带计算所需的部分甚至大部分底层机制，但远不清楚如果这种情况发生，各种特性是否会有一个"大整合"。我们尚未发现有关这种碎片化方法负面影响的任何良好讨论。\footnote{关于该问题的一些初步思考导致了 \cite{sadana2024bringing} 提出利用 Polkadot 技术作为帮助在 roll-up 生态系统之间建立些许兼容性的手段！}

\subsubsection{\textsc{Snark} Roll-ups}

虽然协议的基础并未对 roll-up 的性质做出重大假设，但以太坊的侧带计算策略确实围绕基于 \textsc{snark} 的 roll-up 展开，因此协议正在演化为一种对此有意义的设计。\textsc{Snark} 是一种特殊密码学领域的产物，它允许构建证明以向中立观察者展示执行某个预定义计算的假定结果是正确的。这些证明的验证复杂度趋向于与待证明计算的大小呈亚线性关系，并且不会泄露所述计算的任何内部信息，也不会泄露其可能依赖的任何相关见证数据。

\textsc{Zk-snark} 附带约束。在证明的大小、验证复杂度和生成它的计算复杂度之间存在权衡。非平凡的计算，特别是那种使智能合约如此吸引人的、充满二进制操作的通用计算，很难融入 \textsc{snark} 的模型。

举一个实际的例子，\textsc{risc}-zero（由 \cite{bogli2024assessing} 评估）是一个领先的项目，提供了一个平台来生成由 \textsc{risc-v} 虚拟机执行的计算的 \textsc{snark}——这是一个开源且简洁的 \textsc{risc} 机器架构，受到工具链的良好支持。\cite{koute2024risc0} 最近的基准测试报告表明，与 \textsc{risc}-zero 自己的基准相比，仅证明生成就需要超过 61,000 倍的时间，即使在 32 倍核心数、20,000 倍 \textsc{ram} 和额外的最先进的 \textsc{gpu} 上执行。根据硬件租赁代理商 \url{https://cloud-gpus.com/} 的数据，使用 \textsc{risc}-zero 进行证明的成本乘数是使用 Polka\textsc{vm} 重编译器执行的成本的 66,000,000 倍\footnote{实际上很可能更多，因为这里使用的是消费级设备中的低档"闲置"硬件，而且我们的重编译器未经优化。}。

许多密码学原语变得过于昂贵而无法实际使用，必须用专门的算法和结构来替代。它们通常在其他方面也不是最优的。预期 \textsc{snark} 的使用（如 \cite{cryptoeprint:2019/953} 提出的 \textsc{plonk}），以太坊项目 Dank-sharding 可用性系统的设计使用了以大素数域上多项式承诺为中心的纠删码形式，以使 \textsc{snark} 能够以可接受的性能访问数据的子部分。与本文中的二进制域和 Merklization 等替代方案相比，它导致验证者节点的 \textsc{cpu} 使用率高出数个数量级。

除了基本成本之外，\textsc{snark} 并未提供摆脱去中心化和冗余需求的重大突破，导致了进一步的成本倍增。虽然通过其可验证性质避免了对质押去中心化某些好处的需求，但激励多方做大致相同的工作是确保单方不形成垄断（或多方不形成卡特尔）的要求。证明不正确的状态转换应该是不可能的，然而服务完整性可能以其他方式受到损害；证明生成的暂停——即使只有几分钟——可能对实时金融应用产生重大的经济影响。

去中心化陷阱导致垄断的现实例子确实存在。一个例子就是前述的基于 \textsc{snark} 的交易所框架；虽然名义上服务于去中心化交易所，但实际上是中心化的，Starkware 本身通过生成和提交证明来行使执行交易的垄断权，导致了单点故障——如果 Starkware 的服务受到损害，系统的活性将会受损。

基于 \textsc{snark} 的消除计算信任的策略是否能在成本上与多方加密经济平台竞争，这还有待证明。所有迄今提出的基于 \textsc{snark} 的解决方案都严重依赖加密经济系统来为其构建框架并解决其问题。数据可用性和排序是两个被充分理解为需要加密经济解决方案的领域。

我们注意到 \textsc{snark} 技术正在改进，其背后的密码学家和工程师确实期望在未来几年有所改进。在 \cite{thaler2023technical} 最近的一篇文章中，我们看到一些可信的推测：随着密码学技术的一些最新进展，证明生成的减速可能仅为常规原生执行的 50,000 倍，且其中大部分可以并行化。这比目前的情况要好得多，但仍然比既有的加密经济技术（如 \textsc{Elves}）在成本上竞争所需高出数个数量级。

\subsection{碎片化元网络 (Fragmented Meta-Networks)}

其他项目在通用计算可扩展性方面采取的方向大致集中在两种方法之一；一种可称为\emph{碎片化}方法，另一种可称为\emph{中心化}方法。我们认为两种方法都没有提供令人信服的解决方案。

碎片化方法由 Cosmos（由 \cite{kwon2019cosmos} 提出）和 Avalanche（由 \cite{tanana2019avalanche} 提出）等项目引领。它涉及一个由同质共识机制的网络碎片化而成的系统，但由各自独立激励的验证者集合组成。这与 Polkadot 的单一验证者集合和以太坊宣布的异构 roll-up 策略形成对比——后者部分由在连贯激励框架下运行的同一验证者集合保障。所述碎片化方法的同质性允许相当一致的消息传递机制，有助于向众多连接的网络呈现相当统一的界面。

然而，表面上的一致性是肤浅的。这些网络仅在假设其验证者正确运行时才是无信任的，这些验证者在由经济激励和惩罚最终创造和执行的加密经济安全框架下运行。要以相同的安全级别完成两倍的工作且验证者集合之间没有特殊协调，那么此类系统本质上规定了组建一个具有相同总体激励水平的新网络。

出现了几个问题。首先，存在与 Polkadot 的隔离平行链和以太坊的隔离 roll-up 链类似的缺点：由于持久分片的状态阻止了同步可组合性而缺乏连贯性。

更成问题的是，特别是 Cosmos 提出的碎片化扩展方法不提供同质的安全性——因此也不提供无信任性——保证。网络之间的验证者集合必须被假定为独立选择和激励的，一方在一个网络上的拜占庭行为与另一个网络上采取适当惩罚措施的可能性之间没有因果或概率关系。本质上，这意味着如果验证者合谋损坏或回退一个网络的状态，其影响可能会波及生态系统的其他网络。

这被广泛认为是一个问题，项目提出通过以下两种方式之一来解决。首先，通过从同一验证者集合中抽取来固定跨网络的预期攻击成本（从而固定安全级别）。由 \cite{cosmos2024interchain} 以\emph{复制安全性}名义提出的大规模冗余方式，将要求每个验证者在所有网络上进行验证，并施加相同的激励和惩罚。这在安全提供成本上是不经济的，因为每个网络都需要独立提供与其希望互操作的最安全网络相同水平的激励和惩罚要求。这是为了确保对验证者的经济命题保持不变，并且所有网络的安全命题保持等效。目前，复制安全性并非一个随时可用的无许可服务。我们可以推测，这些惩罚性的经济因素与此有关。

更高效的方法由 OmniLedger 团队 \cite{cryptoeprint:2017/406} 提出，将使验证者非冗余化，在不同网络之间进行分区，并定期、安全和随机地重新分区。与在单一网络上全部验证相比，攻击成本的降低是隐含的，因为存在单个网络意外拥有危及安全数量的恶意验证者的可能性，即使总体上低于该比例。除此之外，它在弱连贯性基础上提供了一种有效的扩展手段。

或者，如 \cite{cryptoeprint:2024/961} 的 \textsc{Elves} 中所述，我们可以利用非冗余分区，将其与验证者参与的提议和审计博弈相结合，以清除和惩罚无效计算，然后要求一个网络的最终性取决于所有因果关联的网络。这是三种方案中最安全和经济最高效的，因为存在一个机制可以高度确信无效转换将在其影响在整个网络生态系统中最终确认之前被识别和纠正。然而，它需要更复杂的逻辑，且其因果关联意味着可添加网络数量的某些上限。

\subsection{高性能全同步网络 (High-Performance Fully Synchronous Networks)}

近年来区块链发展的另一个趋势是通过对数据吞吐量进行"战术性"优化来实现——限制验证者集合的大小或多样性，专注于软件优化，要求验证者之间更高程度的连贯性，对验证者必须拥有的硬件施加苛刻要求，或限制数据可用性。

Solana 区块链由 \cite{yakovenko2018solana} 引入的技术支撑，理论上声称每秒可处理超过 700,000 笔交易，但根据 \cite{ng2024is}，该网络仅被观察到处理了其中的一小部分。底层吞吐量仍然大大超过大多数区块链网络，这归功于有利于最大化同步性能的各种工程优化。结果是一个高度连贯的智能合约环境，其 \textsc{api} 与\emph{黄皮书}以太坊相似（尽管使用不同的底层 \textsc{vm}），但具有近乎即时的包含时间和在包含后即视为最终确认的最终性。

这种方法出现两个问题：首先，将协议定义为高度优化代码库的结果会造成结构性中心化并可能损害韧性。\cite{jha2024solana} 写道"自 2022 年 1 月以来，11 次重大停机事件导致 15 天经历了主要或部分停机"。这在主要区块链中是一个异类，因为绝大多数主要链没有停机时间。停机有各种原因，但通常是由于各子系统中发现的错误。

以太坊，至少在最近之前，提供了最鲜明的对比方案——其经过充分评审的规范、对其加密经济基础的清晰研究以及多个独立实现。也许毫不意外的是，当其部署最广泛的实现中发现一个缺陷并被恶意利用时（如 \cite{hertig2016so} 所述），该网络明显基本未受影响地继续运行。

第二个问题涉及协议在不提供将工作负载分配到单台机器硬件之外的手段时的最终可扩展性。

在主要使用中，历史交易数据和状态都会增长到不切实际的程度。Solana 说明了这可能成为多大的问题。与传统区块链不同，Solana 协议没有为历史数据的归档和后续审查提供解决方案——如果要由第三方从第一性原理证明当前状态是正确的，这至关重要。关于 Solana 如何管理这一点的文献很少，但根据 \cite{solana2023solana}，节点只是将数据放置在由 Google 托管的集中式数据库上。\footnote{较早的节点版本使用了 Arweave 网络——一个去中心化数据存储，但被发现对 Solana 所需的数据吞吐量来说不可靠。}

鼓励 Solana 验证者安装大量的 \textsc{ram} 以帮助在内存中保存其庞大的状态（根据 \cite{solana2024solana}，当前推荐为 512 \textsc{gb}）。没有分而治之的方法，Solana 表明验证者可以合理期望提供的硬件水平决定了完全同步、连贯执行模型的性能上限。硬件要求代表了验证者集合的进入壁垒，不能在不牺牲去中心化和最终透明度的情况下增长。
